% !PRINTED 326
\weberpage{326}
In order to recognise the geometric meaning of the singular points, we imagine the function $f$ ordered by powers of the one variable $x_3$:

\weq{2}{f(x_1, x_2, x_3) = x_3^n f_0 + x_3^{n-1} f_1 + x_3^{n-2} f_2 + \cdots,}
in which $f_0, f_1, f_2, \dots$ are binary forms of the variables $x_1, x_2$, whose degree is indicated by the index, so that $f_0$ is a constant, $f_1$ a linear, $f_2$ a quadratic form, and so on.

Now if a singular point is present, then on account of the invariance of this property we may assume that this point falls at the corner $x_1 = 0$, $x_2 = 0$ of the coordinate triangle, which we shall denote by $\xi_3$. Then $f_0$ and $f_1$ must vanish, and if we then set $x_1 = 0$, the equation $f = 0$ acquires the factor $x_2^2$. This says that two of the points of intersection of the line $x_1 = 0$ coincide in the singular point. But the line $x_2 = 0$ can be any straight line passing through the singular point, and so it follows that the singular point has the property that \gesperrt{every} straight line passing through it cuts the curve there in two coincident points. % [sic: the source prints $x_2=0$ here; the preceding derivation uses the cutting line $x_1=0$]

For this reason the singular points are also called \gesperrt{double points}; by which, however, it is not excluded that more than two points of intersection may also coincide, whereby the \gesperrt{higher singularities} arise.

To the form of the equation (2) we attach a few further simple considerations that will be important in what follows.

If the point $\xi_3$ lies on the curve $f$ but is not a singular point, then the constant $f_0$ must vanish, and $f_1 = 0$ is the equation of the \gesperrt{tangent} of the curve at this point. If we take this tangent for the line $x_1 = 0$, then the equation of the curve, if we assume a constant factor equal to 1, reads:

\weq{3}{f = x_3^{n-1} x_1 + x_3^{n-2} f_2 + \cdots = 0.}
If the quadratic form $f_2$ is divisible by $x_1$, then the equation (3) acquires for $x_1 = 0$ the factor $x_2^3$, and the point $\xi_3$ counts as a threefold point of intersection of $x_1 = 0$ with the curve. Such a point is called a \gesperrt{point of inflection} or \gesperrt{inflexion point} of the curve $f$, and $x_1 = 0$ is called the \gesperrt{inflectional tangent}. The points of inflection do not belong to the singular points.
